\documentclass[10pt,a4paper]{article}
\usepackage[margin=0.85in]{geometry}
\usepackage[T1]{fontenc}
\usepackage[utf8]{inputenc}
\usepackage{lmodern,amsmath,amssymb,graphicx,enumitem,microtype,fancyhdr,xcolor,needspace}
\usepackage[hidelinks]{hyperref}
\definecolor{lightgrey}{gray}{0.94}
\definecolor{midgrey}{gray}{0.35}
\setlength{\parindent}{0pt}
\setlength{\parskip}{5pt}
\setlist[enumerate]{itemsep=5pt,topsep=4pt}
\pagestyle{fancy}\fancyhf{}\renewcommand{\headrulewidth}{0pt}
\fancyfoot[L]{\footnotesize\color{midgrey}Arij Asad}
\fancyfoot[C]{\footnotesize\color{midgrey}Edexcel companion practice}
\fancyfoot[R]{\footnotesize\color{midgrey}\thepage}
\newcounter{question}
\newenvironment{question}[1]{\refstepcounter{question}\Needspace{5\baselineskip}\par\medskip\textbf{\thequestion.\ #1}\par}{\par\vspace{12mm}}
\begin{document}
\hypersetup{pdftitle={Core Pure Mathematics 2 - Original practice and solutions},pdfauthor={Arij Asad}}
\begin{minipage}[c]{0.78\textwidth}{\Large\bfseries Core Pure Mathematics 2}\par Original practice check\end{minipage}\hfill\begin{minipage}[c]{0.17\textwidth}\IfFileExists{PS_logo.png}{\includegraphics[width=\linewidth]{PS_logo.png}}{\rule{0pt}{16mm}}\end{minipage}

\medskip\fcolorbox{black!20}{lightgrey}{\parbox{\dimexpr\linewidth-2\fboxsep-2\fboxrule\relax}{Three short checks on complex roots, hyperbolic functions and differential equations. Give exact answers. Attempt the questions before reading the solutions. These original questions sample selected skills; they are not an official Edexcel paper.}}

\begin{question}{Roots of a complex number}
Find all three complex roots of \(z^3=8\), in the form \(a+bi\).
\end{question}
\begin{question}{Hyperbolic functions}
Solve \(\sinh x=3/4\) for real \(x\), giving an exact answer.
\end{question}
\begin{question}{Second-order differential equations}
Solve \(y^{\prime\prime}-3y^{\prime}+2y=0\), given \(y(0)=1\) and \(y^{\prime}(0)=0\).
\end{question}
\Needspace{7\baselineskip}\textbf{Find the mistake}\par
The curve \(r=2\cos\theta\), for \(-\pi/2\leq\theta\leq\pi/2\), traces a circle once. A student calculates its enclosed area as \(\int_{-\pi/2}^{\pi/2}2\cos\theta\,\mathrm{d}\theta=4\). Identify the error and find the area.

\vfill\small Further practice: \href{https://maths.arijasad.com/tmua-paper-archive.html}{TMUA papers and solutions} and \href{https://maths.arijasad.com/maths-topic-practice.html}{maths topic guides}.
\newpage
{\Large\bfseries Hints and worked solutions}\par\medskip
\Needspace{8\baselineskip}\subsection*{1. Roots of a complex number}
\textit{Hint: Write \(8=8e^{2k\pi i}\), then take cube roots with three distinct arguments.}\par
The roots have modulus \(8^{1/3}=2\) and arguments \(2k\pi/3\), for \(k=0,1,2\).

Using \(z=2(\cos(2k\pi/3)+i\sin(2k\pi/3))\) gives \(2\), \(-1+i\sqrt3\) and \(-1-i\sqrt3\). These are distinct and account for all roots of the cubic.

\textbf{Answer:} \(2,\ -1+i\sqrt3,\ -1-i\sqrt3\).

\Needspace{8\baselineskip}\subsection*{2. Hyperbolic functions}
\textit{Hint: Use the exponential definition of \(\sinh x\), and set \(u=e^x>0\).}\par
\((e^x-e^{-x})/2=3/4\). With \(u=e^x>0\), this becomes \(u-1/u=3/2\).

Multiplying by \(2u\) gives \(2u^2-3u-2=0\), or \((2u+1)(u-2)=0\).

The root \(u=-1/2\) is impossible because \(e^x>0\). Thus \(u=2\), so \(x=\ln2\).

\textbf{Answer:} \(x=\ln2\).

\Needspace{8\baselineskip}\subsection*{3. Second-order differential equations}
\textit{Hint: Find the two auxiliary roots, then use both initial conditions to determine the constants.}\par
The auxiliary equation is \(m^2-3m+2=(m-1)(m-2)=0\), giving distinct roots \(1\) and \(2\).

The general solution is \(y=Ae^x+Be^{2x}\), with \(y^{\prime}=Ae^x+2Be^{2x}\).

The initial conditions give \(A+B=1\) and \(A+2B=0\). Subtraction gives \(B=-1\), then \(A=2\). Hence \(y=2e^x-e^{2x}\).

\textbf{Answer:} \(y=2e^x-e^{2x}\).

\Needspace{9\baselineskip}\subsection*{Find the mistake: correction}
The student has integrated \(r\); the correct area formula is \(\tfrac12\int r^2\,\mathrm{d}\theta\).

Here the area is \[\begin{aligned}A&=\frac12\int_{-\pi/2}^{\pi/2}4\cos^2\theta\,\mathrm{d}\theta\\&=\int_{-\pi/2}^{\pi/2}(1+\cos2\theta)\,\mathrm{d}\theta\\&=\left[\theta+\frac12\sin2\theta\right]_{-\pi/2}^{\pi/2}\\&=\pi.\end{aligned}\]

As a check, multiplying the polar equation by \(r\) gives \(x^2+y^2=2x\), or \((x-1)^2+y^2=1\): a circle of radius \(1\), whose area is \(\pi\).

\Needspace{6\baselineskip}\subsection*{What to practise next}\begin{enumerate}
\item For complex roots, organise modulus and argument separately and check that you have found the required number of distinct roots.
\item Derive unfamiliar hyperbolic identities from exponentials instead of transferring trigonometric signs from memory.
\item Check differential-equation solutions twice: substitute into the equation and then verify every initial or boundary condition.
\end{enumerate}\end{document}
